\section{Robustness and Implementation}
\label{sec:robustness}

The purpose of the robustness analysis is not to search retrospectively for the calibration with the highest historical performance. It is to determine whether the conclusions reported in the previous section survive economically plausible changes in implementation, risk estimation, exposure limits, sample composition, and investor preferences.

The two empirical channels require different tests. Dynamic equity--bond allocation is primarily exposed to turnover, model instability, and the cost of responding to changing estimated state probabilities. The model-based direct variance-payoff approximation is additionally exposed to monthly roll costs, negatively skewed payoff shocks, errors in lagged risk estimation, notional constraints, and the concentration of returns in particular market episodes.

All sensitivity comparisons preserve the timing rules described in the methodology. Direct-variance parameter grids are aligned to a common sample within each sensitivity dimension. Consequently, a 12-month and a 60-month risk-estimation specification are compared only over dates available to both. Appendix B summarizes the parameter grid and supplementary robustness diagnostics, while the full machine-readable grids are retained under \texttt{outputs/tables/}; the main text focuses on the results required to assess the economic conclusion.

\subsection{Allocation implementation robustness}

The principal allocation models require substantially more trading than the static benchmarks. In the United States allocation sample, average turnover is 1.47\% for 60/40 and 1.53\% for equal weighting, compared with 25.35\% for HMM RV + Log VRP and 36.69\% for RSM RV + Raw VRP. In Europe, the strongest dynamic specifications similarly require turnover between approximately 15\% and 34\%, compared with less than 2\% for the balanced portfolios.

This implementation burden matters because the gross advantage of the dynamic models is small. Table~\ref{tab:allocation_cost_robustness} shows the transaction-cost sensitivity of the selected United States HMM.

\begin{table}[htbp]
\centering
\small
\caption{United States allocation Sharpe ratios under transaction costs}
\label{tab:allocation_cost_robustness}
\begin{tabular}{lrrr}
\toprule
Cost
& 60/40
& 1/N
& HMM RV + Log VRP \\
\midrule
0 bps
& 1.026
& 1.028
& 1.055 \\
10 bps
& 1.024
& 1.025
& 1.019 \\
25 bps
& 1.021
& 1.022
& 0.963 \\
50 bps
& 1.016
& 1.016
& 0.869 \\
\bottomrule
\end{tabular}
\end{table}

Before transaction costs, the HMM produces a higher Sharpe ratio than the two balanced portfolios. At the 10-basis-point baseline, the ranking reverses. At 25 and 50 basis points, the HMM deteriorates more rapidly because its turnover is much larger. Its drawdown advantage is more persistent than its Sharpe advantage: at baseline costs, maximum drawdown remains -16.67\%, compared with -20.06\% for 60/40 and -19.10\% for equal weighting. The appropriate interpretation is therefore defensive rather than one of broad performance dominance.

A no-trade band produces only a modest improvement. Increasing the band from zero to 10\% raises the selected HMM Sharpe ratio from approximately 1.019 to 1.026 and lowers average turnover from approximately 25.30\% to 23.53\%. Small adjustments are removed, but large changes in estimated state probabilities still require substantial trading.

Partial rebalancing is more effective. Instead of moving fully to the model target each month, the implemented portfolio moves only a fraction of the distance. Table~\ref{tab:partial_rebalancing} compares full rebalancing with an adjustment speed of 0.25.

\begin{table}[htbp]
\centering
\small
\caption{Partial rebalancing of the selected United States HMM}
\label{tab:partial_rebalancing}
\begin{tabularx}{\textwidth}{Xrrrrr}
\toprule
Implementation
& Return
& Vol.
& Sharpe
& MDD
& Turnover \\
\midrule
Full rebalancing
& 8.36\%
& 8.24\%
& 1.019
& -16.67\%
& 25.35\% \\
Partial, $\lambda=0.25$
& 8.69\%
& 8.62\%
& 1.014
& -18.04\%
& 9.10\% \\
\bottomrule
\end{tabularx}
\end{table}

The 0.25 implementation reduces turnover by approximately two thirds while changing the Sharpe ratio only marginally. Maximum drawdown remains below that of 60/40 and buy-and-hold equity. This makes the HMM operationally more plausible, but it does not make it a statistically or economically dominant allocation strategy.

The European allocation conclusion is even less sensitive to implementation choices. The strongest HMM reaches a Sharpe ratio of 0.218, the strongest RSM 0.281, and the strongest machine-learning strategy 0.274. These values remain close to or below the simple benchmarks while requiring materially greater turnover. There is therefore no strong baseline allocation advantage for implementation refinements to preserve.

\subsection{Direct-variance transaction-cost sensitivity}

The direct-variance cost treatment differs from portfolio rebalancing. Each one-month modeled variance exposure settles and a new exposure must be initiated. Costs are therefore charged on the full absolute monthly notional rather than only on the change in notional.

The tested cost rates are 0, 10, 25, and 50 basis points per unit of monthly notional. Table~\ref{tab:direct_costs} reports the complete selected-strategy cost grid.

\begin{table}[htbp]
\centering
\scriptsize
\caption{Direct-variance transaction-cost sensitivity}
\label{tab:direct_costs}
\begin{tabularx}{\textwidth}{Xrrrrrrr}
\toprule
Market / cost
& Return
& Vol.
& Sharpe
& MDD
& MV CEQ
& vs 60/40
& vs 1/N \\
\midrule
US 0 bps
& 6.26\%
& 7.18\%
& 0.885
& -23.42\%
& 5.06\%
& -0.69\%
& -0.43\% \\
US 10 bps
& 6.23\%
& 7.17\%
& 0.882
& -23.46\%
& 5.04\%
& -0.71\%
& -0.45\% \\
US 25 bps
& 6.19\%
& 7.17\%
& 0.877
& -23.53\%
& 5.00\%
& -0.74\%
& -0.48\% \\
US 50 bps
& 6.13\%
& 7.16\%
& 0.869
& -23.63\%
& 4.95\%
& -0.80\%
& -0.54\% \\
EU 0 bps
& 13.14\%
& 9.70\%
& 1.329
& -21.35\%
& 10.54\%
& 9.32\%
& 9.04\% \\
EU 10 bps
& 13.08\%
& 9.70\%
& 1.324
& -21.36\%
& 10.49\%
& 9.27\%
& 8.99\% \\
EU 25 bps
& 12.99\%
& 9.69\%
& 1.316
& -21.38\%
& 10.41\%
& 9.19\%
& 8.91\% \\
EU 50 bps
& 12.84\%
& 9.69\%
& 1.303
& -21.42\%
& 10.28\%
& 9.06\%
& 8.78\% \\
\bottomrule
\end{tabularx}
\end{table}

The selected United States strategy is not highly sensitive to the tested cost range because its average active notional is low. This does not improve the welfare conclusion: its mean--variance certainty-equivalent difference is already negative at zero costs and becomes slightly more negative as costs rise.

The European result is more informative. Even at 50 basis points, the Sharpe ratio remains 1.303 and the annualized mean--variance certainty-equivalent advantage remains 9.06 percentage points relative to 60/40 and 8.78 percentage points relative to equal weighting. The result is therefore not generated by assuming zero or near-zero stylized roll costs. This test does not establish that real derivative implementation costs would lie inside the tested range.

\subsection{Sensitivity to the risk-estimation window}

The baseline payoff-volatility estimate uses a 36-month lookback. Robustness tests consider 12, 24, 36, and 60 months. These specifications are evaluated on the common sample required by the longest window, so the 36-month values in Table~\ref{tab:lookback_robustness} differ slightly from the unrestricted baseline.

\begin{table}[htbp]
\centering
\scriptsize
\caption{Direct-variance risk-estimation-window sensitivity}
\label{tab:lookback_robustness}
\begin{tabularx}{\textwidth}{Xrrrrrrr}
\toprule
Market / window
& Return
& Vol.
& Sharpe
& MDD
& MV CEQ
& vs 60/40
& vs 1/N \\
\midrule
US 12m
& 8.10\%
& 16.15\%
& 0.578
& -50.69\%
& 2.81\%
& -3.30\%
& -2.94\% \\
US 24m
& 6.38\%
& 9.94\%
& 0.678
& -32.59\%
& 4.27\%
& -1.85\%
& -1.49\% \\
US 36m
& 6.50\%
& 7.10\%
& 0.926
& -23.46\%
& 5.32\%
& -0.80\%
& -0.44\% \\
US 60m
& 6.26\%
& 6.87\%
& 0.922
& -23.00\%
& 5.15\%
& -0.97\%
& -0.60\% \\
EU 12m
& 18.35\%
& 11.21\%
& 1.570
& -15.92\%
& 14.45\%
& 11.74\%
& 11.78\% \\
EU 24m
& 15.66\%
& 10.46\%
& 1.453
& -19.97\%
& 12.46\%
& 9.75\%
& 9.79\% \\
EU 36m
& 14.15\%
& 9.79\%
& 1.409
& -21.36\%
& 11.40\%
& 8.69\%
& 8.73\% \\
EU 60m
& 13.92\%
& 8.93\%
& 1.513
& -17.21\%
& 11.52\%
& 8.80\%
& 8.84\% \\
\bottomrule
\end{tabularx}
\end{table}

The United States result is sensitive to short risk-estimation windows. The 12-month specification realizes volatility of 16.15\%, a maximum drawdown of -50.69\%, and a Sharpe ratio of only 0.578. A short history can therefore assign excessive exposure before a large payoff shock. The 36- and 60-month specifications are substantially more stable, although their welfare differences remain negative relative to both benchmarks.

The European conclusion is stable across the entire grid. Every tested window produces a Sharpe ratio above 1.40 on the aligned sample and a large positive certainty-equivalent difference. The strongest ex post calibration is not the object of the test; the important result is that the European conclusion does not depend on a unique volatility lookback.

\subsection{Notional-cap sensitivity}

The maximum absolute notionals tested are 5\%, 10\%, 15\%, 25\%, and 50\%. Caps above 15\% are effectively non-binding for the selected United States strategy, while caps above 10\% are non-binding for the selected European strategy. Table~\ref{tab:cap_robustness} therefore reports the economically informative subset of the grid.

\begin{table}[htbp]
\centering
\small
\caption{Selected notional-cap sensitivity}
\label{tab:cap_robustness}
\begin{tabularx}{\textwidth}{Xrrrrrr}
\toprule
Market / cap
& Return
& Vol.
& Sharpe
& MDD
& MV CEQ
& vs benchmarks \\
\midrule
US 5\%
& 5.37\%
& 6.53\%
& 0.836
& -23.46\%
& 4.39\%
& negative \\
US 10\%
& 6.25\%
& 7.17\%
& 0.884
& -23.46\%
& 5.05\%
& negative \\
US 25\%
& 6.23\%
& 7.17\%
& 0.882
& -23.46\%
& 5.04\%
& negative \\
EU 5\%
& 12.07\%
& 7.94\%
& 1.482
& -17.95\%
& 10.19\%
& positive \\
EU 10\%
& 13.08\%
& 9.70\%
& 1.324
& -21.36\%
& 10.49\%
& positive \\
EU 25\%
& 13.08\%
& 9.70\%
& 1.324
& -21.36\%
& 10.49\%
& positive \\
\bottomrule
\end{tabularx}
\end{table}

The 5\% European cap is particularly informative. Despite reducing permitted exposure by 80\% relative to the baseline 25\% cap, the strategy still produces a 12.07\% annualized return, 7.94\% volatility, a Sharpe ratio of 1.482, and a maximum drawdown of -17.95\%. Its mean--variance certainty-equivalent advantages remain 8.97 percentage points relative to 60/40 and 8.69 percentage points relative to equal weighting. The European result is therefore not driven by occasional reliance on the maximum baseline notional.

In the United States, changing the cap does not reverse the welfare conclusion. Reducing the cap to 5\% lowers both return and volatility, but the strategy remains below the traditional benchmarks in mean--variance certainty-equivalent terms.

\subsection{Subperiod stability}

Time stability is essential because an attractive full-sample average can conceal a result concentrated in a small part of the history. Table~\ref{tab:subperiod_robustness} reports the selected strategies across the principal sample partitions.

\begin{table}[htbp]
\centering
\scriptsize
\caption{Subperiod stability of the selected direct-variance strategies}
\label{tab:subperiod_robustness}
\begin{tabularx}{\textwidth}{Xrrrrr}
\toprule
Market / period
& Return
& Vol.
& Sharpe
& MDD
& Delta CEQ range \\
\midrule
US full
& 6.23\%
& 7.17\%
& 0.882
& -23.46\%
& -0.71\% / -0.45\% \\
US first half
& 6.64\%
& 5.84\%
& 1.131
& -6.11\%
& 1.68\% / 1.43\% \\
US second half
& 5.83\%
& 8.32\%
& 0.726
& -23.46\%
& -3.11\% / -2.33\% \\
US pre-Covid
& 5.40\%
& 7.19\%
& 0.770
& -13.60\%
& -1.13\% / -1.15\% \\
US Covid+
& 7.87\%
& 7.17\%
& 1.096
& -15.37\%
& 0.13\% / 0.92\% \\
EU full
& 13.08\%
& 9.70\%
& 1.324
& -21.36\%
& 9.27\% / 8.99\% \\
EU first half
& 12.77\%
& 7.07\%
& 1.743
& -6.33\%
& 11.09\% / 10.16\% \\
EU second half
& 13.38\%
& 11.79\%
& 1.132
& -21.36\%
& 7.44\% / 7.81\% \\
EU pre-Covid
& 12.81\%
& 7.30\%
& 1.695
& -6.33\%
& 10.43\% / 9.66\% \\
EU Covid+
& 13.39\%
& 11.92\%
& 1.121
& -21.16\%
& 7.92\% / 8.22\% \\
\bottomrule
\end{tabularx}
\end{table}

The United States result is clearly time dependent. The first half produces positive welfare differences, while the second half produces materially negative differences. Performance is also stronger from the Covid period onward than in the preceding sample. These changes prevent a claim of stable benchmark dominance.

The European strategy remains economically strong in every tested subperiod. Volatility and drawdown increase in the second half and after Covid, but annualized return remains close to 13\% and the mean--variance certainty-equivalent advantage remains above seven percentage points against both benchmarks even in the weaker partitions.

\subsection{Crisis exclusions}

Crisis exclusions remove the Global Financial Crisis, the European sovereign-debt crisis, Volmageddon, the Covid shock, and the 2022 inflation and interest-rate shock individually and jointly. The test asks whether the full-sample conclusion is mechanically generated by a few volatility episodes.

The United States result remains dependent on which stress window is removed. Some exclusions improve Sharpe ratio and drawdown, but they do not uniformly reverse the welfare conclusion. This reinforces the interpretation of a strategy whose historical value depends materially on the timing of a limited number of volatility shocks.

The European result is substantially more stable. Table~\ref{tab:crisis_robustness} compares the full sample with the joint exclusion of all designated major-crisis observations.

\begin{table}[htbp]
\centering
\small
\caption{European crisis-exclusion robustness}
\label{tab:crisis_robustness}
\begin{tabularx}{\textwidth}{Xrrrrrr}
\toprule
Scenario
& Obs.
& Return
& Vol.
& Sharpe
& MDD
& MV CEQ \\
\midrule
Full sample
& 170
& 13.08\%
& 9.70\%
& 1.324
& -21.36\%
& 10.49\% \\
Exclude all major crises
& 141
& 17.38\%
& 8.80\%
& 1.877
& -14.05\%
& 14.59\% \\
\bottomrule
\end{tabularx}
\end{table}

The joint exclusion removes 29 observations. The Sharpe ratio rises to 1.877 and maximum drawdown falls to -14.05\%. Mean--variance certainty-equivalent advantages remain approximately 8.57 percentage points relative to 60/40 and 8.90 percentage points relative to equal weighting. The European result is therefore not created solely by gains concentrated in the designated crisis windows.

This test does not imply that trading through crisis periods would be frictionless. It only establishes that the historical statistical result survives when those observations are removed.

\subsection{Risk-aversion sensitivity}

The final robustness dimension varies the investor risk-aversion coefficient over $\gamma \in \{1,3,5,10\}$. Table~\ref{tab:gamma_robustness} reports the mean--variance certainty-equivalent differences.

\begin{table}[htbp]
\centering
\small
\caption{Risk-aversion sensitivity}
\label{tab:gamma_robustness}
\begin{tabular}{lrrr}
\toprule
Market
& $\gamma$
& vs 60/40
& vs 1/N \\
\midrule
US
& 1
& -1.65\%
& -0.91\% \\
US
& 3
& -1.18\%
& -0.68\% \\
US
& 5
& -0.71\%
& -0.45\% \\
US
& 10
& 0.47\%
& 0.12\% \\
EU
& 1
& 8.81\%
& 9.10\% \\
EU
& 3
& 9.04\%
& 9.05\% \\
EU
& 5
& 9.27\%
& 8.99\% \\
EU
& 10
& 9.83\%
& 8.84\% \\
\bottomrule
\end{tabular}
\end{table}

The United States strategy is below both benchmarks for risk-aversion coefficients from one to five. It becomes slightly favorable under the mean--variance criterion at $\gamma=10$, where its lower volatility receives greater weight. This isolated result is not sufficient to establish general superiority.

The European certainty-equivalent advantage remains positive and economically large for every tested value of $\gamma$. CRRA certainty equivalents decline more rapidly at high risk aversion because the negatively skewed payoff distribution receives a larger penalty, but the central welfare conclusion remains unchanged.

\subsection{Robustness synthesis and implementation boundary}

The robustness evidence sharpens the cross-market interpretation.

For allocation, implementation refinements improve practicality more than they improve economic ranking. Partial rebalancing materially reduces turnover in the selected United States HMM, but neither the United States nor European allocation evidence becomes a robust case of benchmark dominance.

For the United States direct-variance strategy, moderate transaction costs are not the principal weakness. The more important issues are sensitivity to short risk-estimation windows, substantial differences across subperiods, and the absence of consistent certainty-equivalent superiority. The evidence supports a defensive variance-carry interpretation, not a universal return premium.

For Europe, the selected VRP-positive strategy survives every principal robustness dimension tested: costs up to 50 basis points, volatility lookbacks from 12 to 60 months, a maximum notional reduced to 5\%, separate first and second halves, pre- and post-Covid samples, joint exclusion of major crises, and risk-aversion coefficients from one to ten.

This robustness remains conditional on the modeling framework. The direct-variance layer still uses a volatility-index-based strike proxy, a trailing realized-variance settlement proxy, stylized monthly costs, and no explicit collateral, margin, market-impact, dealer-spread, or intramonth mark-to-market process. Robustness within the approximation does not convert the result into an observed or exactly replicable variance-swap return.

The appropriate conclusion is therefore that the European direct variance-payoff result is robust within the tested model-based framework, whereas the United States evidence is substantially more conditional and the dynamic allocation evidence remains economically close to simple benchmarks.
